\documentclass[10pt,a4paper]{amsart}
\usepackage[margin=19mm]{geometry}
\usepackage{amsmath,amssymb,mathtools}
\usepackage[scr=rsfs,cal=euler]{mathalfa}
\usepackage{xfrac}
\usepackage[unicode,hidelinks,bookmarks=false]{hyperref}
\newcommand{\ie}{i.e.\ }
\newcommand{\cf}{cf.\ }
\newcommand{\resp}{respectively\ }
\newcommand{\Subsection}{\subsection}
\providecommand{\nobreakdash}{\nobreak\hbox{-}}
\setlength{\emergencystretch}{2em}
\multlinegap0pt
\usepackage{bm}
\usepackage{bbm}
\usepackage{dsfont}
\usepackage{xspace}
\usepackage{cancel}
\usepackage{mathtools}
\usepackage{amsthm}
\makeatletter
\@ifundefined{theorem}{\newtheorem{theorem}{Theorem}}{}
\@ifundefined{lemma}{\newtheorem{lemma}[theorem]{Lemma}}{}
\makeatother
\newcommand{\N}{\mathbb{N}}
\newcommand{\cB}{\mathcal{B}}
\newcommand{\cF}{\mathcal{F}}
\title[ADDQ: Adaptive Distributional Double Q-Learning]{ADDQ: Adaptive Distributional Double Q-Learning}
\author{Leif D\"oring, Benedikt Wille, Maximilian Birr, Mihail B\^irsan, Martin Slowik}
\date{Source: arXiv:2506.19478v1 (CC BY 4.0)}
\begin{document}
\maketitle

\noindent\emph{Source and licence.} ADDQ: Adaptive Distributional Double Q-Learning. Version arXiv:2506.19478v1 (CC BY 4.0), distributed under the Creative Commons Attribution 4.0 International licence, as recorded in the arXiv licence metadata for this exact version and shown on the abstract page https://arxiv.org/abs/2506.19478v1. Full rights notice: licenses/arxiv-2506.19478-CC-BY-4.0.txt.\par\smallskip

\noindent\textit{Editorial scope.} Counted unit: the Lemma whose source label is \texttt{lem:RM} in the article (source0.tex lines 1683--1686, source label \texttt{lem:RM}), delivered together with the article's own complete original proof of it (source0.tex lines 1687--1715, one \texttt{proof} environment of 29 lines). No mathematics is written, repaired or re-derived: statement and proof are the article's own text line for line. Required context retained uncounted: none. Every definition, display and label of those lines is retained unchanged. Omitted from this package and not redistributed: 1--1682, 1716--1970. Nothing inside a retained excerpt is omitted or altered: every retained body is delivered exactly as the article's lines are written, so no omission receipt of any kind accompanies this package. Citations: the retained excerpts carry no citation command at all, so no bibliography entry of the article is transcribed and no reference is redistributed. Reference adaptation: none was needed and none was made; every reference inside the delivered unit resolves inside it, so no unresolved reference or citation is produced. Printed slips kept literal: no printed word, symbol, hypothesis or number of the counted statement or of its proof was altered. Layout and class: the article's own class is \texttt{article} with packages including microtype, graphicx, subfigure, booktabs, amsmath, floatflt, adjustbox, float, hyperref, icml2025, amssymb, mathtools, amsthm, cleveref; the wrapper is the lane's pinned \texttt{amsart} editorial wrapper at 10pt on A4, which restates the theorem-like environments the retained text needs and adds the article's own macro definitions that the retained text needs (\texttt{\textbackslash N}, \texttt{\textbackslash cB}, \texttt{\textbackslash cF}), with extra wrapper packages (bm, bbm, dsfont, xspace, cancel, mathtools). The delivered numbering is the unit's own. Assets: no retained span contains a figure environment, a graphics-inclusion command, a PDF-page inclusion, a table or a TikZ picture; no image file is copied into this package, the package declares no asset dependency and no image file is redistributed with it. Licence: version arXiv:2506.19478v1 of this article is distributed under the Creative Commons Attribution 4.0 International licence, as recorded in the arXiv licence metadata for this exact version and shown on the abstract page https://arxiv.org/abs/2506.19478v1. A full rights notice is delivered at licenses/arxiv-2506.19478-CC-BY-4.0.txt. Characters: every retained excerpt is pure ASCII, so the delivered engine, XeLaTeX with its default Latin Modern text font, reports no missing character.
\begin{lemma}\label{lem:RM}
    Let $(\alpha_t)_{t\in\N_0}$ be a sequence fulfilling the Robbins-Monro conditions and $(Y_t)_{t\in\N}$ an $iid$ sequence of Bernoulli(0.5) random variables, i.e. $\mathbb P(Y_t=1)=\mathbb P(Y_t=0)=0.5$ for all $t\in\N_0$.
    Then $(\alpha_tY_t)_{t\in\N_0}$ also fulfills the Robbins-Monro condition.
\end{lemma}

\begin{proof}
    The almost sure convergence of the summed squares is obviously fulfilled due to 
    \[\sum_{t=0}^\infty(\alpha_tY_t)^2\leq\sum_{t=0}^\infty \alpha_t^2<\infty\quad \text{almost surely.}\]
    Due to independence of each $Y_t$ with $\{Y_n|n\in\N_0,\ n\neq t\}$ as well as with $\alpha=(\alpha_t)_{t=0}^\infty$ we will consider a two stage experiment, where we first draw the sequence $\alpha=(\alpha_t)_{t=0}^\infty$ and then independently of this realization sample the $iid$ sequence $Y=(Y_t)_{t=0}^\infty$. Due to the independence the joint measure of $\alpha$ and $Y$ is the product measure. Consider the product space $(\Omega,\cF,\mathbb P)=(\Omega_\alpha\times\Omega_Y,\cF_\alpha\otimes\cF_Y,\mathbb P^{\otimes\N}_\alpha\otimes\mathbb P^{\otimes\N}_Y)$ where $\Omega_\alpha,\Omega_Y=[0,1]^\N$, $\cF_\alpha,\cF_Y=\cB([0,1])^{\otimes \N}$ . Then, using that $\sum_{t=0}^\infty\alpha_t=\infty$ $\mathbb P_\alpha$-almost surely, we have
    \begin{align*}   
    \mathbb P\Big(\sum_{t=0}^\infty\alpha_tY_t=\infty\Big)
    =&\int_{\Omega_\alpha}
    \mathbb P_Y\Big(\sum_{t=0}^\infty\alpha_tY_t=\infty\Big)
    d\mathbb P_\alpha(\alpha)\\
            =&\int_{\{(\alpha_t)_{t=0}^\infty\in \Omega_\alpha:\sum_{t=0}^\infty\alpha_t=\infty\}}\mathbb P_Y\Big(\sum_{t=0}^\infty\alpha_tY_t=\infty\Big)d\mathbb P_\alpha(\alpha)\\
            \overset{(a)}{=}&\int_{\{(\alpha_t)_{t=0}^\infty\in \Omega_\alpha:\sum_{t=0}^\infty\alpha_t=\infty\}}1\ d\mathbb P_\alpha(\alpha)\\
            =& 1,
    \end{align*} 
    where $(a)$ can be seen as follows. Consider any deterministic sequence $(b_t)\subseteq[0,1]$ fulfilling $\sum_{t=0}^\infty b_t=\infty$. Then
    \[\infty= \sum_{t=0}^\infty b_t=\sum_{t=0}^\infty b_tY_t+\sum_{t=0}^\infty b_t\textbf{1}_{Y_t=0}.\]
    Now notice that $A=\sum_{t=0}^\infty b_tY_t$ and $B=\sum_{t=0}^\infty b_t\textbf{1}_{Y_t=0}$ are identically distributed and since the sum of $A$ and $B$ is always infinity, almost surely either one of them is infinite. Given the identical distribution, we infer
    \[\mathbb P_Y(\sum_{t=0}^\infty b_tY_t=\infty)>0.\]
    But since $(b_tY_t)$ is an independent sequence of random variables and the event that the infinite sum diverges is in the tail sigma algebra, the Kolmogorov 0-1 law yields:
    \[\mathbb P_Y(\sum_{t=0}^\infty b_tY_t=\infty)=1.\]
    %Now, for the first condition we recursively construct a deterministic sequence as follows. For $K\in\N$, define
    %\[t_K:=\min\{\n\in\N | n>t_{K-1}\text{ and }\sum_{t=0}^\infty\alpha_t>K \text{ almost surely}\}.\]
    %Since $\sum_{t=0}^\infty\alpha_t=\infty $ almost surely, such a sequence exists. Further define a deterministic sequence $(b_t)$ that is one exactly at the %indices $t_K$, for $K\in\N$, and $0$ otherwise. This is
    %\[b_t=\begin{cases}
    %    1\quad ,\ if\ \exists K\in\N,\ s.t.\ t=t_K\\
    %    0\quad otherwise.
    %\end{cases}\]
    %Then, we have
    %\[\sum_{t=0}^\infty\alpha_tY_t \]
\end{proof}

\end{document}
